



\subsection{Sweep Rule}
The sweep-rule is the core of our construction. In a decoding cycle, each of the vertices reads the syndrome on its positive $X$-checks (checks that are associated with the positive $d^{\prime}-1$ dimensional faces of the vertex) and, based on what it sees, applies a Pauli correction. The sweep-rule defines the correction to be over the positive qubits such that it will turn off the positive syndrome. It's not hard to see that for the $(2,2)$-Toric complex, the Toom's and the sweep rules are the same. Compared to other correction cycles on topological codes, the main advantage of the sweep is that it is independent of time. The invention of the sweep-rule is attributed to Aleksander Kubica and John Preskill \cite{Kubica_2019}. Here we extend their work to cubical complexes.

\begin{definition}[Sweep Rule]
  \label{def:sweeprule}
  Let $z$ be an assignment on the $d^{\prime}$-faces. The \textit{sweep rule decoding procedure} at a vertex $v \in \mathcal{G}_{n}^{d}$ is defined as follows:
  
    Measure the checks rooted at $v$ and look for a rooted assignment $\tilde{z}$ at $v$ such that $(\partial^{-}\tilde{z})|_{v+} = \partial^{-}(z)|_{v+}$. In other words, the $X$-checks rooted at $v$ have the same syndromes for $z$ and $\tilde{z}$. If such an assignment exists, flip $\tilde{z}$.

\end{definition}
A full decoding cycle is composed of first applying the sweep-rule regarding the $X$-syndrome, then moving into the dual code (and reordering the qubits) by $\mathbf{H}_{\phi}$, and then applying the sweep-rule again with the change that the positive direction is flipped to be the negative direction. At the end, we return to the original code by applying $\mathbf{H}_{\phi}$.

\begin{definition}[Sweep cycle]
We define the \textit{sweep cycle} to be the following procedure:
\begin{enumerate}
    \item For each vertex, perform the sweep rule.
    \item Apply the Hadamard transform $\mathbf{H}^{\otimes n}$ followed by the $\phi$-transpose, and flip the positive direction.
    \item Again, apply the sweep rule at each vertex.
    \item Reverse stage 2 by reordering the qubits according to the $\phi$-transpose, and then apply the Hadamard transform $\mathbf{H}^{\otimes n}$. Flip the positive direction again.
\end{enumerate}
\end{definition}


To have our constructions yield circuits that have $\infty$-extensions \Cref{def:extension}, the sweep-rule cycles have to be 'decoding gadgets' \Cref{def:decgate}.  
\begin{claim}
  The local gates that perform the sweep rule are decoding gadgets.
\end{claim}
\begin{proof}
  A local sweep update acts on a constant-size neighborhood around a vertex, uses only local syndrome information, and may be implemented with fresh ancillas that are reset after the update. Any dependence on faults inside the gadget can therefore be pushed onto the data wires as an effective Pauli operator supported on the encoded register, while the reset ancillas carry no residual information. This is exactly the form required in the definition of a decoding gadget.
\end{proof}


